% Luc Destombe --- licence CC BY-NC-SA 4.0
% https://creativecommons.org/licenses/by-nc-sa/4.0/deed.fr
% Fichier autonome : compiler avec pdflatex, deux fois (signets, nombre de pages).
\documentclass[11pt,a4paper]{article}
\makeatletter
\newif\ifeleve
\elevefalse
\elevetrue

\newif\iflycee@palatino
\lycee@palatinotrue

\RequirePackage[utf8]{inputenc}
\RequirePackage[T1]{fontenc}
\RequirePackage[french]{babel}
\RequirePackage{etoolbox}

\newcommand{\ShorthandsOff}{\@ifundefined{shorthandoff}{}{\shorthandoff{;:!?}}}
\newcommand{\ShorthandsOn}{\@ifundefined{shorthandon}{}{\shorthandon{;:!?}}}
\ShorthandsOff
\AtBeginDocument{\ShorthandsOn}
\AtBeginEnvironment{tikzpicture}{\ShorthandsOff}

\RequirePackage{geometry}
\RequirePackage{fancyhdr}
\RequirePackage{lastpage}
\RequirePackage{multicol}
\RequirePackage{array}
\RequirePackage{enumitem}
\RequirePackage{xcolor}
\RequirePackage{needspace}

\RequirePackage{amsmath,amssymb}
\RequirePackage{mathpazo}
\DeclareMathSizes{10.95}{10.95}{8}{5.5}
\begingroup\catcode`\_=\active \catcode`\^=\active
\gdef_#1{\sb{\scriptscriptstyle #1}}
\gdef^#1{\sp{\scriptscriptstyle\lycee@moinsepais #1}}
\endgroup
\newcommand\lycee@moins{\mathbin{%
  \setbox\z@\hbox{$\m@th\scriptscriptstyle\lycee@moinsorig$}%
  \dimen@=.55\wd\z@ \advance\dimen@-.5pt
  \hbox to.75\wd\z@{\hss$\m@th\scriptscriptstyle\vcenter{\hbox{%
    \tikz\draw[line width=.5pt,line cap=round](0,0)--(\the\dimen@,0);}}$\hss}}}
\begingroup\lccode`\~=`\- \lowercase{\endgroup\let~}\lycee@moins
\newcommand\lycee@moinsepais{\mathcode`\-="8000 }
\AtBeginDocument{\mathchardef\lycee@moinsorig=\mathcode`\-\relax}
\AtBeginDocument{\catcode`\_=12 \mathcode`\_="8000
  \catcode`\^=12 \mathcode`\^="8000 }
\newcommand\lycee@exposants{%
  \fontdimen13\textfont2=.48\fontdimen6\textfont2
  \fontdimen14\textfont2=.48\fontdimen6\textfont2
  \fontdimen15\textfont2=.48\fontdimen6\textfont2
  \fontdimen13\scriptfont2=.48\fontdimen6\scriptfont2
  \fontdimen14\scriptfont2=.48\fontdimen6\scriptfont2
  \fontdimen15\scriptfont2=.48\fontdimen6\scriptfont2}
\every@math@size\expandafter{\the\every@math@size\lycee@exposants}
\newlength{\retraitcalculs}
\setlength{\retraitcalculs}{2em}

\RequirePackage{tikz}
\usetikzlibrary{arrows.meta,calc,babel,shapes.misc}
\RequirePackage[most]{tcolorbox}
\tcbuselibrary{skins,breakable}

\definecolor{coulDef}{HTML}{1F4E79}
\definecolor{coulProp}{HTML}{7030A0}
\definecolor{coulRem}{HTML}{C55A11}
\definecolor{coulEx}{HTML}{2E7D32}
\definecolor{coulExo}{HTML}{B03A2E}
\definecolor{coulPart}{HTML}{1F4E79}
\definecolor{coulMeth}{HTML}{1B6E4F}
\definecolor{coulCourbe}{HTML}{0B5ED7}
\definecolor{coulCourbeB}{HTML}{C00000}
\definecolor{coulCourbeC}{HTML}{2E7D32}
\definecolor{coulGrille}{HTML}{8C8C8C}
\definecolor{coulRep}{HTML}{1B6E4F}
\definecolor{coulBar}{HTML}{2E7D32}

\newcommand{\R}{\mathbb{R}}
\newcommand{\Rs}{\mathbb{R}^{*}}
\renewcommand{\le}{\leqslant}
\renewcommand{\ge}{\geqslant}
\newcommand{\vect}[1]{\overrightarrow{#1}}
\newcommand{\norme}[1]{\left\lVert #1 \right\rVert}
\newcommand{\intervalle}[2]{\left[#1\,;#2\right]}
\RequirePackage{cancel}
\renewcommand{\CancelColor}{\color{coulExo}}
\newcommand{\fois}[1]{\times{\color{coulExo}#1}}
\newcommand{\barre}[1]{\cancel{#1}}

\tikzset{grille/.style={coulGrille,line width=0.4pt}}
\newcommand{\quadrillage}[4]{\draw[grille,step=1] (#1,#2) grid (#3,#4);}
\tikzset{
  croix/.style={cross out,draw,line width=0.6pt,inner sep=0pt,minimum size=4pt},
  croixdroite/.style={croix,rotate=45,line width=0.8pt},
}
\newcommand{\pt}[3]{\path (#1) node[croix]{} node[#2,font=\small,inner sep=2.5pt]{$#3$};}

\newcommand{\lycee@repere}[4]{%
  \draw[grille,step=1] (#1,#2) grid (#3,#4);
  \draw[-{Stealth[length=2mm]},thick] (#1,0)--({#3+0.4},0) node[below left=-1pt]{$x$};
  \draw[-{Stealth[length=2mm]},thick] (0,#2)--(0,{#4+0.4}) node[below left=-1pt]{$y$};
  \node[below left,font=\scriptsize,inner sep=1.5pt] at (0,0) {$0$};}
\newcommand{\lycee@gradx}[1]{\foreach \k/\l in {#1}{%
  \draw (\k,0.07)--(\k,-0.07) node[below,font=\scriptsize,inner sep=1.5pt]{$\l$};}}
\newcommand{\lycee@grady}[1]{\foreach \k/\l in {#1}{%
  \draw (0.07,\k)--(-0.07,\k) node[left,font=\scriptsize,inner sep=1.5pt]{$\l$};}}
\newcommand{\lycee@intff}[2]{\left[#1\,;#2\right]}
\newcommand{\lycee@intfo}[2]{\left[#1\,;#2\right[}
\newcommand{\lycee@intof}[2]{\left]#1\,;#2\right]}
\newcommand{\lycee@intoo}[2]{\left]#1\,;#2\right[}
\newcommand{\lycee@ens}[1]{\left\{#1\right\}}
\AtBeginDocument{%
  \providecommand{\repere}{\lycee@repere}%
  \providecommand{\gradx}{\lycee@gradx}%
  \providecommand{\grady}{\lycee@grady}%
  \providecommand{\Cf}{\mathcal{C}\sb{\scriptscriptstyle f}}%
  \providecommand{\Cg}{\mathcal{C}\sb{\scriptscriptstyle g}}%
  \providecommand{\intff}{\lycee@intff}%
  \providecommand{\intfo}{\lycee@intfo}%
  \providecommand{\intof}{\lycee@intof}%
  \providecommand{\intoo}{\lycee@intoo}%
  \providecommand{\ens}{\lycee@ens}}

\newlist{questions}{enumerate}{3}
\setlist[questions,1]{label=\arabic*\textdegree),leftmargin=*,itemsep=2pt,topsep=3pt}
\setlist[questions,2]{label=\alph*),leftmargin=*,itemsep=1pt,topsep=2pt}
\setlist[questions,3]{label=\roman*),leftmargin=*,itemsep=1pt,topsep=1pt}
\setlist[itemize]{label=\textbullet,leftmargin=*,itemsep=2pt,topsep=3pt}

\newcommand{\besoinplace}[1]{\par\penalty\@M\begingroup
  \dimen@=#1\relax\dimen@ii=\pagegoal\advance\dimen@ii-\pagetotal
  \ifdim\dimen@>\dimen@ii\ifdim\dimen@ii>\z@\vfil\fi\break\fi\endgroup}

\newcommand{\lycee@pied}{\thepage/\pageref{LastPage}}
\newcommand{\entete}[2]{%
  \pagestyle{fancy}\fancyhf{}%
  \fancyhead[L]{\footnotesize\color{coulPart}\textbf{#1}}%
  \fancyhead[R]{\footnotesize\color{coulPart}\textbf{#2}}%
  \fancyfoot[C]{\footnotesize\lycee@pied}%
  \renewcommand{\headrulewidth}{0.4pt}%
  \renewcommand{\headrule}{\hbox to\headwidth{%
    \color{coulPart!40}\leaders\hrule height \headrulewidth\hfill}}}

\RequirePackage{ccicons}
\newcommand{\lycee@licence}{{\scriptsize\color{gray}%
  \copyright\ Luc Destombe --- \ccbyncsa\ CC BY-NC-SA 4.0}}
\newif\iflycee@licence \lycee@licencetrue
\newcommand{\sanslicence}{\lycee@licencefalse}
\AtBeginDocument{\iflycee@licence\AddToHookNext{shipout/foreground}{%
  \put(0,0){\rlap{\kern\dimexpr1in+\hoffset+\oddsidemargin+\textwidth\relax
    \llap{\raisebox{-\dimexpr1in+\voffset+\topmargin+\headheight+\headsep
      +\textheight+\footskip\relax}{\lycee@licence}}}}}\fi}

\newcommand{\formulesserrees}{%
  \setlength{\abovedisplayskip}{3pt}\setlength{\belowdisplayskip}{3pt}%
  \setlength{\abovedisplayshortskip}{2pt}\setlength{\belowdisplayshortskip}{3pt}}

\setlength{\parindent}{0pt}

\RequirePackage[implicit=false,hidelinks,pdfencoding=auto,psdextra,bookmarksopen,
  bookmarksopenlevel=1,pdfcreator={},pdfproducer={}]{hyperref}
\RequirePackage{bookmark}
\hypersetup{pdfauthor={Luc Destombe},pdfkeywords={CC BY-NC-SA 4.0}}
\newcounter{lycee@signet}
\newcommand{\signet}[3][]{\stepcounter{lycee@signet}%
  \edef\lycee@dest{\if\relax\detokenize{#1}\relax
    signet.\arabic{lycee@signet}\else#1\fi}%
  \hypertarget{\lycee@dest}{}\bookmark[level=#2,dest=\lycee@dest]{#3}}
\pdfstringdefDisableCommands{%
  \def\R{R}\def\vect#1{#1}\def\Cf{Cf}\def\Cg{Cg}\def\fois#1{×#1}%
  \def\barre#1{#1}\def\intervalle#1#2{[#1;#2]}\def\intff#1#2{[#1;#2]}%
  \def\textsuperscript#1{#1}\def\dfrac#1#2{#1/#2}\def\frac#1#2{#1/#2}%
  \def\vec#1{vecteur #1}}
\geometry{top=2cm,bottom=1cm,left=1cm,right=1cm,headheight=14pt,footskip=0.6cm}

\RequirePackage{pgfplots}
\pgfplotsset{compat=1.18}
\RequirePackage{tkz-tab}

\newcounter{partiecours}
\renewcommand{\thepartiecours}{\Roman{partiecours}}
\newcommand{\partiecours}[1]{%
  \refstepcounter{partiecours}%
  \Needspace*{8\baselineskip}%
  \bigskip
  \noindent\begin{tcolorbox}[enhanced,colback=coulPart,colframe=coulPart,
      boxrule=0pt,arc=2pt,left=8pt,right=8pt,top=4pt,bottom=4pt,
      before skip=6pt,after skip=10pt]
    \color{white}\bfseries\large \thepartiecours{} --- #1\signet{1}{\thepartiecours{} --- #1}
  \end{tcolorbox}%
  \addcontentsline{toc}{section}{\thepartiecours{} --- #1}%
}

\newtcolorbox{boitecours}[3][]{%
  enhanced,breakable,
  colback=white,colframe=#2,coltitle=white,
  fonttitle=\bfseries,
  attach boxed title to top left={xshift=8pt,yshift=-\tcboxedtitleheight/2},
  boxed title style={colback=#2,arc=2pt,boxrule=0pt},
  boxrule=0.9pt,arc=2pt,
  left=8pt,right=8pt,top=8pt,bottom=6pt,
  title={#3},#1}

\newcounter{methode}
\newenvironment{definitionbox}[1][Définition]%
  {\begin{boitecours}[phantom={\signet{2}{#1}}]{coulDef}{#1}}{\end{boitecours}}
\newenvironment{proprietebox}[1][Propriété]%
  {\begin{boitecours}[phantom={\signet{2}{#1}}]{coulProp}{#1}}{\end{boitecours}}
\newenvironment{methodebox}[1][Méthode]{\stepcounter{methode}%
  \begin{boitecours}[phantom={\signet[methode-\arabic{methode}]{2}{#1}}]{coulMeth}{#1}}%
  {\end{boitecours}}

\newcommand{\etape}[1]{\par\smallskip\textbf{\color{coulMeth}#1}\par\nobreak\smallskip}
\newcommand{\enonce}[1]{\emph{#1}\par\smallskip}
\newcommand{\reponse}[1]{\par\smallskip
  {\footnotesize\itshape\color{black!55}Réponses : #1}}
\newcommand{\demo}[1]{\textbf{\color{coulMeth}Démonstration.} #1}
\newcommand{\afaire}[1]{%
  \par\smallskip
  \noindent{\small\color{coulExo}$\blacktriangleright$\ \textbf{À faire :}
    fiche d'exercices du chapitre, #1.}\par\smallskip}

\newenvironment{calculs}{\setlength{\jot}{5pt}\@fleqntrue\@mathmargin\retraitcalculs
  \setlength{\abovedisplayskip}{4pt}\setlength{\belowdisplayskip}{4pt}%
  \setlength{\abovedisplayshortskip}{4pt}\setlength{\belowdisplayshortskip}{4pt}%
  \csname alignat*\endcsname{2}}%
  {\csname endalignat*\endcsname}
\newcommand{\rem}[1]{\text{\small\itshape\color{black!60}#1}}
\newcommand{\explique}[2]{\par\smallskip\noindent
  \begin{minipage}[t]{0.57\linewidth}\raggedright #1\end{minipage}\hfill
  \begin{minipage}[t]{0.40\linewidth}\raggedright\leavevmode
    {\small\itshape\color{black!60}#2}\end{minipage}\par}
\newcommand{\cas}[1]{\par\smallskip\noindent\textbf{\color{coulExo}#1}\nobreak}
\newcommand{\calc}[1]{\par\smallskip\textbf{\color{coulExo}#1}\par\nobreak}

\newtcolorbox{remarquebox}[1][Remarque]{%
  enhanced,breakable,colback=white,colframe=coulRem,
  boxrule=0pt,leftrule=2.5pt,arc=0pt,outer arc=0pt,
  left=8pt,right=6pt,top=5pt,bottom=5pt,
  before upper={\textbf{\color{coulRem}#1}\par\smallskip}}

\newtcolorbox{exemplebox}[1][Exemple]{%
  enhanced,breakable,colback=white,colframe=coulEx,
  boxrule=0pt,leftrule=2.5pt,arc=0pt,outer arc=0pt,
  left=8pt,right=6pt,top=5pt,bottom=5pt,
  before upper={\textbf{\color{coulEx}#1}\par\smallskip}}

\pgfplotsset{
  repere/.style={
    axis lines=middle,
    axis line style={-{Stealth[length=2.4mm]},thick},
    every axis x label/.style={at={(ticklabel* cs:1.0)},anchor=west},
    every axis y label/.style={at={(ticklabel* cs:1.0)},anchor=south},
    label style={font=\footnotesize},
    tick label style={font=\scriptsize},
    grid=both,
    major grid style={grille},
    minor tick num=0,
    tick align=inside,
    clip mode=individual,
  },
}
\tikzset{
  courbe/.style={coulCourbe,line width=1.1pt,smooth},
  courbeB/.style={coulCourbeB,line width=1.1pt,smooth},
  courbeC/.style={coulCourbeC,line width=1.1pt,smooth},
}

\newlist{etapes}{enumerate}{1}
\setlist[etapes]{label=\textbf{\arabic*.},leftmargin=*,itemsep=1pt,topsep=2pt}

\newlist{expressions}{enumerate}{1}
\setlist[expressions]{label={\Alph*\ $=$},leftmargin=*,itemsep=3pt,topsep=3pt}

\newcounter{exo}
\renewcommand{\theexo}{\ifnum\value{exo}<10 0\fi\arabic{exo}}
\newtcolorbox{exercice}[1][]{%
  enhanced,breakable,
  colback=white,colframe=coulExo,
  boxrule=0.9pt,arc=2pt,
  left=8pt,right=8pt,top=8pt,bottom=6pt,
  coltitle=white,fonttitle=\bfseries,
  attach boxed title to top left={xshift=8pt,yshift=-\tcboxedtitleheight/2},
  boxed title style={colback=coulExo,arc=2pt,boxrule=0pt},
  phantom={\signet{2}{Exercice \arabic{exo}}},
  title={Exercice \theexo\ifx\relax#1\relax\else{} --- #1\fi}}
\newenvironment{exo}[1][\relax]{\refstepcounter{exo}\begin{exercice}[#1]}{\end{exercice}}

  \RequirePackage{comment}
  \excludecomment{correction}
  \renewcommand{\reponse}[1]{}
  \newcommand{\autableau}{\hfill{\footnotesize\itshape\color{black!45}(au tableau)}}
  \newcommand{\etapeexemples}{%
    \par\smallskip\textbf{\color{coulMeth}Exemples}\autableau\par\nobreak\smallskip}
  \newcommand{\etapetableau}[1]{%
    \par\smallskip\textbf{\color{coulMeth}#1}\autableau\par\nobreak\smallskip}
  \newtcolorbox{exempletableau}[1][Exemple]{%
    enhanced,breakable,colback=white,colframe=coulEx,
    boxrule=0pt,leftrule=2.5pt,arc=0pt,outer arc=0pt,
    left=8pt,right=6pt,top=4pt,bottom=4pt,
    before skip=5pt,after skip=5pt,
    before upper={\textbf{\color{coulEx}#1}\autableau\par\smallskip}}

\newcommand{\coursEntete}{}
\newcommand{\coursNiveau}{}
\newcommand{\coursTitre}{}
\newcommand{\coursSurTitre}{}
\newcommand{\coursSousTitre}{}

\newcommand{\enteteCours}[1]{\renewcommand{\coursEntete}{#1}}
\newcommand{\chapitrecours}[2]{\enteteCours{Chapitre #1 --- #2}}
\newcommand{\niveaucours}[1]{\renewcommand{\coursNiveau}{#1}}
\newcommand{\titrecours}[3][]{%
  \renewcommand{\coursTitre}{#2}%
  \renewcommand{\coursSurTitre}{#1}%
  \renewcommand{\coursSousTitre}{#3}}

\pagestyle{fancy}
\fancyhf{}
\fancyhead[L]{\footnotesize\color{coulPart}\textbf{\coursEntete}}
\fancyhead[R]{\footnotesize\color{coulPart}\textbf{\coursNiveau}}
\fancyfoot[C]{\footnotesize\thepage}
\renewcommand{\headrulewidth}{0.4pt}
\renewcommand{\headrule}{\hbox to\headwidth{\color{coulPart!40}\leaders\hrule height \headrulewidth\hfill}}

\setlength{\parskip}{2pt}

\AtBeginDocument{%
  \begin{center}
    \begin{tcolorbox}[enhanced,width=0.72\linewidth,colback=white,colframe=coulPart,
        boxrule=1.6pt,arc=3pt,halign=center,top=10pt,bottom=10pt]
      \ifx\coursSurTitre\empty
        {\Huge\bfseries\color{coulPart} \coursTitre}\\[4pt]
      \else
        {\Huge\bfseries\color{coulPart} \coursTitre}\\[3pt]
        {\Large\bfseries\color{coulPart} \coursSurTitre}\\[5pt]
      \fi
      {\normalsize \coursSousTitre}%
      \\[2pt]{\small\itshape Version élève}
    \end{tcolorbox}
  \end{center}%
}
\makeatother

\chapitrecours{1}{Calcul littéral}
\niveaucours{1\textsuperscript{re} spécialité mathématiques}
\titrecours{CALCUL LITTÉRAL}{Cours --- Première spécialité}

\begin{document}

\begin{remarquebox}[Comment utiliser ce cours]
Chaque partie commence par ce qu'il faut \textbf{savoir} (définitions, théorèmes), puis
vient ce qu'il faut \textbf{savoir faire} : une méthode, avec des
\textbf{exemples} faits en classe, puis des exercices
\og\textbf{À vous de jouer}\fg{} à chercher seul.

\smallskip
Sur cette feuille, seuls les \textbf{énoncés} des exemples figurent : la résolution,
signalée par la mention \emph{(au tableau)}, est à rédiger dans le cahier.

\end{remarquebox}

\partiecours{Réduire, développer, identités remarquables}

\begin{definitionbox}[Définition 1 --- Réduire]
Des termes \textbf{semblables} ont la même partie littérale : $3x^{2}$ et $-5x^{2}$ sont
semblables, $3x^{2}$ et $3x$ ne le sont pas.

\smallskip
\textbf{Réduire une somme}, c'est regrouper les termes semblables en additionnant leurs
coefficients. \textbf{Réduire un produit}, c'est multiplier les nombres entre eux et les
lettres entre elles.

\smallskip
On peut \textbf{toujours} réduire un produit : $3x\times 2x=6x^{2}$. On ne peut
\textbf{pas toujours} réduire une somme : $3x^{2}+2x$ n'a pas de termes semblables, elle
reste ainsi.
\end{definitionbox}

\begin{methodebox}[Méthode 1 --- Réduire une somme ou un produit]
\etapeexemples
Réduire et ordonner les expressions suivantes.
\begin{questions}[label=\alph*)]
  \item $A=5x^{2}-3x+7-2x^{2}+8x-4$
  \item $B=2x\times\left(-4x^{2}\right)\times 3$
\end{questions}

\etape{À vous de jouer}
Réduire :
\[
  C=6x-4+x^{2}-9x+10 \qquad D=2x^{2}+5x-3x^{2}+x \qquad
  E=-3x\times 5x \qquad F=7x\times\left(-2x^{2}\right)
\]

\end{methodebox}

\begin{definitionbox}[Définition 2 --- Développer, factoriser]
\begin{minipage}[t]{0.46\linewidth}
\vspace{2pt}
\textbf{Développer}, c'est transformer un produit en une somme.

\smallskip
\textbf{Factoriser}, c'est transformer une somme en un produit.
\end{minipage}\hfill
\begin{minipage}[t]{0.50\linewidth}
\vspace{0pt}
\centering
\begin{tikzpicture}[>={Stealth[length=3mm]}]
  \node (g) at (0,0) {$3x(x-2)$};
  \node (d) at (3.6,0) {$3x^{2}-6x$};
  \node[font=\scriptsize,anchor=north] (gl) at (g.south) {forme factorisée};
  \node[font=\scriptsize,anchor=north] (dl) at (d.south) {forme développée};
  \draw[->,line width=1pt,coulEx] (g.north east) to[bend left=32]
        node[above,font=\scriptsize\bfseries,coulEx]{DÉVELOPPER} (d.north west);
  \draw[->,line width=1pt,coulExo] (dl.south west) to[bend left=32]
        node[below,font=\scriptsize\bfseries,coulExo]{FACTORISER} (gl.south east);
\end{tikzpicture}
\end{minipage}
\end{definitionbox}

\begin{proprietebox}[Propriété 1 --- Distributivité et identités remarquables]
Pour tous réels $a$, $b$, $c$, $d$ et $k$ :
\begin{center}
$\begin{aligned}
  k(a+b) &= ka+kb & (a+b)^{2} &= a^{2}+2ab+b^{2}\\
  (a+b)(c+d) &= ac+ad+bc+bd \qquad & (a-b)^{2} &= a^{2}-2ab+b^{2}\\
  & & (a+b)(a-b) &= a^{2}-b^{2}
\end{aligned}$
\end{center}
\end{proprietebox}

\begin{methodebox}[Méthode 2 --- Développer et réduire]
\etapeexemples
Développer et réduire les expressions suivantes.
\begin{questions}[label=\alph*)]
  \item $A=(2x+5)(x-3)$
  \item $B=(3x-2)^{2}$
  \item $C=(x+4)^{2}-(x+1)(x-1)$
\end{questions}

\etape{À vous de jouer}
Développer et réduire :
\[
  D=(x-6)(4x+1) \qquad E=(5x+2)^{2} \qquad F=(2x-7)(2x+7) \qquad G=3x(x-2)-(x-5)^{2}
\]

\end{methodebox}

\afaire{exercices 1 à 3}

\partiecours{Substituer une valeur}

\begin{definitionbox}[Définition 3 --- Valeur d'une expression, valeur interdite]
Calculer la \textbf{valeur} d'une expression pour $x=a$, c'est remplacer chaque $x$ par
$a$, puis effectuer les calculs.

\smallskip
On ne peut pas diviser par $0$ : une valeur qui annule un dénominateur est une
\textbf{valeur interdite}. L'expression n'a pas de valeur pour ce nombre.
\end{definitionbox}

\begin{methodebox}[Méthode 3 --- Calculer une image]
\etapeexemples
\begin{questions}[label=\alph*)]
  \item $P(x)=2x^{2}-3x+1$ : calculer $P(-2)$.
  \item $f(x)=\dfrac{x+5}{x-1}$ : donner la valeur interdite, puis calculer $f(3)$ et
        $f\!\left(\dfrac{1}{3}\right)$.
  \item Avec la même fonction $f$, calculer $f\!\left(\sqrt{3}\right)$ sans racine au
        dénominateur.
\end{questions}

\etape{À vous de jouer}
\begin{questions}
  \item Avec $Q(x)=x^{2}-4x+3$, calculer $Q(-1)$ et $Q\!\left(\sqrt{2}\right)$.
  \item Avec $g(x)=\dfrac{x+4}{x-2}$, donner la valeur interdite, puis calculer $g(4)$ et
        $g\!\left(\dfrac{1}{2}\right)$.
  \item Calculer $g\!\left(\sqrt{5}\right)$, sans racine au dénominateur.
\end{questions}

\end{methodebox}

\afaire{exercices 4 à 6}

\partiecours{Factoriser une expression}

Pour factoriser, on dispose de deux outils : le \textbf{facteur commun},
$ka+kb=k(a+b)$, et les \textbf{identités remarquables} lues de droite à gauche,
$a^{2}+2ab+b^{2}=(a+b)^{2}$, $a^{2}-2ab+b^{2}=(a-b)^{2}$ et $a^{2}-b^{2}=(a-b)(a+b)$.

\begin{methodebox}[Méthode 4 --- Factoriser]
\etapeexemples
Factoriser les expressions suivantes.
\begin{questions}[label=\alph*)]
  \item $A=6x^{2}+10x$
  \item $B=(x+3)(2x-1)-(x+3)(x+4)$
  \item $C=9x^{2}-30x+25$
  \item $D=(x+2)^{2}-49$
\end{questions}

\etape{À vous de jouer}
Factoriser :
\[
  E=12x^{2}-8x \qquad F=(2x+1)(x-3)+(2x+1)(4x+5) \qquad G=x^{2}+14x+49
\]
\[
  H=16-(x-1)^{2} \qquad I=25x^{2}-4
\]

\end{methodebox}

\afaire{exercices 7 à 9}

\partiecours{Équations}

\begin{proprietebox}[Propriété 2 --- Produit nul]
Un produit est nul si et seulement si l'un au moins de ses facteurs est nul :
\[
  A\times B=0 \iff A=0 \;\text{ ou }\; B=0
\]
\end{proprietebox}

\begin{proprietebox}[Propriété 3 --- Équation $x^{2}=k$]
\begin{itemize}
  \item Si $k>0$, l'équation $x^{2}=k$ a deux solutions : $-\sqrt{k}$ et $\sqrt{k}$.
  \item Si $k=0$, elle a une seule solution : $0$.
  \item Si $k<0$, elle n'a pas de solution, car un carré n'est jamais négatif.
\end{itemize}
\end{proprietebox}

\begin{methodebox}[Méthode 5 --- Résoudre une équation]
\etapeexemples
Résoudre dans $\R$ les équations suivantes.
\begin{questions}[label=\alph*)]
  \item $5x-3=2x+9$
  \item $(2x-6)(5x+1)=0$
  \item $x^{2}=36$ ; $x^{2}=11$ ; $x^{2}=-9$ ; $(x+1)^{2}=25$
  \item $x^{2}=3x$
\end{questions}

\etape{À vous de jouer}
Résoudre dans $\R$ :
\[
  4(x-2)=x+7 \qquad (3x+12)(7-x)=0 \qquad x^{2}=49
\]
\[
  x^{2}=-5 \qquad (x-2)^{2}=16 \qquad x^{2}=5x
\]

\end{methodebox}

\begin{methodebox}[Méthode 6 --- Mettre un problème en équation]
\etapeexemples
Un samedi, une salle d'escalade a vendu $120$ entrées : des entrées adultes à $14$~€ et
des entrées jeunes à $9$~€. La recette est de $1\,430$~€. Combien d'entrées de chaque
sorte ont été vendues ?

\etape{À vous de jouer}
\begin{questions}
  \item Une association a vendu $300$ billets pour un concert : en prévente à $18$~€, ou
        sur place à $25$~€. La recette est de $6\,100$~€. Combien de billets ont été
        vendus en prévente ?
  \item La longueur d'un rectangle dépasse sa largeur de $4$~cm, et son périmètre mesure
        $36$~cm. Quelles sont ses dimensions ?
\end{questions}

\end{methodebox}

\afaire{exercices 10 à 12}

\partiecours{Fractions rationnelles}

\begin{definitionbox}[Définition 4 --- Fraction rationnelle]
Une \textbf{fraction rationnelle} est un quotient de deux expressions en $x$, comme
$\dfrac{3x+1}{x-2}$. Ses valeurs interdites sont celles qui annulent le dénominateur.

\smallskip
L'\textbf{ensemble de définition} est l'ensemble des réels qui ne sont pas interdits :
ici $D=\R\setminus\{2\}$, qui se lit « $\R$ privé de $2$ ».
\end{definitionbox}

\begin{methodebox}[Méthode 7 --- Trouver les valeurs interdites, simplifier]
\etapeexemples
\begin{questions}[label=\alph*)]
  \item Donner l'ensemble de définition de $f(x)=\dfrac{4}{x+3}$ et de
        $g(x)=\dfrac{x-1}{2x-5}$.
  \item Simplifier $h(x)=\dfrac{x^{2}-4}{x^{2}+2x}$, en précisant les valeurs interdites.
\end{questions}

\etape{À vous de jouer}
\begin{questions}
  \item Donner l'ensemble de définition de $k(x)=\dfrac{7}{x-6}$ et de
        $m(x)=\dfrac{3x+1}{4x+2}$.
  \item Simplifier $n(x)=\dfrac{x^{2}-25}{x^{2}-5x}$, en précisant les valeurs interdites.
\end{questions}

\end{methodebox}

\begin{methodebox}[Méthode 8 --- Écrire sous la forme d'un seul quotient]
\etapeexemples
Écrire sous la forme d'un seul quotient, en précisant les valeurs interdites.
\begin{questions}[label=\alph*)]
  \item $A=2+\dfrac{3}{x-1}$
  \item $B=\dfrac{1}{x}-\dfrac{2}{x+3}$
\end{questions}

\etape{À vous de jouer}
Écrire sous la forme d'un seul quotient, en précisant les valeurs interdites :
\[
  C=5-\dfrac{2}{x+4} \qquad D=\dfrac{3}{x-2}+\dfrac{1}{x} \qquad E=\dfrac{x}{x+1}-2
\]

\end{methodebox}

\afaire{exercices 13 à 16}

\partiecours{Équations quotient}

\begin{proprietebox}[Propriété 4 --- Quotient nul]
Un quotient est nul si et seulement si son numérateur est nul, son dénominateur ne
l'étant pas :
\[
  \dfrac{A}{B}=0 \iff A=0 \;\text{ et }\; B\ne 0
\]
\end{proprietebox}

\begin{methodebox}[Méthode 9 --- Résoudre une équation quotient]
\etapeexemples
Résoudre dans $\R$ les équations suivantes.
\begin{questions}[label=\alph*)]
  \item Un quotient égal à $0$ : $\dfrac{2x-8}{x+5}=0$.
  \item Un quotient égal à un nombre : $\dfrac{x+7}{x-2}=4$.
  \item Un quotient égal à un quotient : $\dfrac{x+1}{x-3}=\dfrac{x+2}{x-1}$.
\end{questions}

\begin{remarquebox}[Deux remarques]
\begin{itemize}
  \item Une solution du numérateur peut être interdite : on l'écarte. Ainsi
        $\dfrac{x^{2}-9}{x-3}=0$ donne $x=3$ ou $x=-3$, mais $3$ est interdit :
        $S=\{-3\}$.
  \item Le produit en croix donne aussi le résultat en c), mais il ne sert plus pour les
        inéquations (partie VII) : on retient la méthode du quotient nul.
\end{itemize}
\end{remarquebox}

\etape{À vous de jouer}
Résoudre dans $\R$, en donnant d'abord les valeurs interdites :
\[
  \dfrac{3x+9}{x-4}=0 \qquad \dfrac{x+9}{x+1}=3 \qquad \dfrac{x+4}{x+2}=\dfrac{x-1}{x-2}
\]

\end{methodebox}

\afaire{exercices 17 à 19}

\partiecours{Inéquations et tableaux de signes}

\begin{proprietebox}[Propriété 5 --- Signe de $ax+b$ (avec $a\ne 0$)]
$ax+b$ s'annule pour $x=-\dfrac{b}{a}$ ; \textbf{après} cette valeur, il a le signe de
$a$, \textbf{avant}, le signe contraire.
\begin{center}
\begin{minipage}{0.47\linewidth}\centering
{\small $a>0$}\\[2pt]
\begin{tikzpicture}[scale=0.9]
  \tkzTabInit[lgt=1.6,espcl=2.2]{$x$/1,$ax+b$/1}{$-\infty$,$-\frac{b}{a}$,$+\infty$}
  \tkzTabLine{,-,z,+,}
\end{tikzpicture}
\end{minipage}\hfill
\begin{minipage}{0.47\linewidth}\centering
{\small $a<0$}\\[2pt]
\begin{tikzpicture}[scale=0.9]
  \tkzTabInit[lgt=1.6,espcl=2.2]{$x$/1,$ax+b$/1}{$-\infty$,$-\frac{b}{a}$,$+\infty$}
  \tkzTabLine{,+,z,-,}
\end{tikzpicture}
\end{minipage}
\end{center}
\end{proprietebox}

\begin{exemplebox}[Exemple --- signe de $2x-6$ et de $-4x+2$]
$2x-6$ s'annule en $3$ et $a=2>0$ : négatif avant $3$, positif après.

\smallskip
$-4x+2$ s'annule en $\dfrac{1}{2}$ et $a=-4<0$ : positif avant $\dfrac{1}{2}$, négatif
après.
\end{exemplebox}

\begin{methodebox}[Méthode 10 --- Résoudre une inéquation]
\etapeexemples
Résoudre dans $\R$ les inéquations suivantes.
\begin{questions}[label=\alph*)]
  \item Premier degré : $2(x+3)\ge 5x-9$.
  \item Quotient comparé à $0$ : $\dfrac{x-2}{x+4}<0$.
  \item Quotient comparé à un nombre : $\dfrac{6}{x+2}<2$.
\end{questions}

\etape{À vous de jouer}
Résoudre dans $\R$ :
\[
  3(1-x)<x+11 \qquad \dfrac{x+3}{x-1}\ge 0 \qquad \dfrac{4}{x-3}\ge 2
\]

\end{methodebox}

\afaire{exercices 20 à 22}

\partiecours{Composition de polynômes}

\begin{definitionbox}[Définition 5 --- Fonction polynôme]
Une \textbf{fonction polynôme} s'écrit $P(x)=a_{n}x^{n}+\dots+a_{1}x+a_{0}$, où les
\textbf{coefficients} $a_{n}$, \dots, $a_{0}$ sont des réels. Si $a_{n}\ne 0$, le nombre
$n$ est le \textbf{degré} de $P$.

\smallskip
Exemple : $P(x)=4x^{3}-x+7$ est de degré $3$.
\end{definitionbox}

\begin{definitionbox}[Définition 6 --- Composée]
\textbf{Composer} deux fonctions, c'est mettre une fonction à l'intérieur d'une autre.

\smallskip
La \textbf{composée} de $P$ par $Q$ est la fonction $P\circ Q$ définie par
$(P\circ Q)(x)=P\bigl(Q(x)\bigr)$ : dans $P\bigl(Q(x)\bigr)$, on met la fonction $Q$ à
l'intérieur de la fonction $P$.

\smallskip
Concrètement, on remplace chaque $x$ de $P(x)$ par $Q(x)$.
\end{definitionbox}

\begin{methodebox}[Méthode 11 --- Calculer une composée]
\etapeexemples
Soient $P(x)=x^{2}-3x$ et $Q(x)=x+2$. Calculer $(P\circ Q)(x)$, puis $(Q\circ P)(x)$.

\etape{À vous de jouer}
\begin{questions}
  \item $P(x)=2x^{2}+1$ et $Q(x)=x-3$ : calculer $(P\circ Q)(x)$ et $(Q\circ P)(x)$.
  \item $P(x)=x^{2}+4x$ et $Q(x)=3x-1$ : calculer $(P\circ Q)(x)$, puis $(P\circ Q)(1)$
        de deux façons.
\end{questions}

\end{methodebox}

\afaire{exercices 23 à 25}

\end{document}
